\subsubsection{Giải chi tiết Bài tập a (Phương pháp đơn hình)}

\textbf{1. Đưa bài toán về dạng chuẩn tắc}

Vì tất cả các ràng buộc chính đều là bất phương trình dạng "$\le$", ta cộng thêm vào vế trái của mỗi ràng buộc một biến phụ không âm $x_4, x_5, x_6 \ge 0$. 
Ta có mô hình bài toán dạng chuẩn như sau:

\begin{alignat*}{4}
    f(x) = \quad & 3x_1 & {}+{} & 2x_2 & {}+{} & 5x_3 & {}+{} 0x_4 + 0x_5 + 0x_6 \rightarrow \max \\
    \text{s.t.} \quad & x_1 & {}+{} & x_2 & {}+{} & 2x_3 & {}+{} x_4 = 4000 \\
    & x_1 & {}+{} & x_2 & {}+{} & 3x_3 & {}+{} x_5 = 9000 \\
    & 3x_1 & {}+{} & x_2 & {}+{} & 2x_3 & {}+{} x_6 = 8000 \\
    & \rlap{$x_j \ge 0, \quad j = \overline{1,6}$}
\end{alignat*}

Hệ vector điều kiện ứng với các biến phụ $x_4, x_5, x_6$ tạo thành ma trận đơn vị $I_3$, do đó ta chọn luôn $(x_4, x_5, x_6)$ làm cơ sở xuất phát.

\vspace{0.3cm}
\textbf{2. Quá trình giải bằng Bảng đơn hình}

\textbf{Bảng 1 (Bước lặp 0):}
\begin{table}[H]
    \centering
    \renewcommand{\arraystretch}{1.3}
    \begin{tabular}{|c|c|c|c|c|c|c|c|c|c|}
        \hline
        \multirow{2}{*}{\textbf{Cơ sở}} & \multirow{2}{*}{$\mathbf{C_b}$} & \multirow{2}{*}{\textbf{PA}} & $x_1$ & $x_2$ & $x_3$ & $x_4$ & $x_5$ & $x_6$ & \multirow{2}{*}{\textbf{Tỷ số $\lambda$}} \\
        \cline{4-9}
        & & & 3 & 2 & 5 & 0 & 0 & 0 & \\
        \hline
        $x_4$ & 0 & 4000 & 1 & 1 & \textbf{2} & 1 & 0 & 0 & $4000/2 = \mathbf{2000} \rightarrow$ \\
        \hline
        $x_5$ & 0 & 9000 & 1 & 1 & 3 & 0 & 1 & 0 & $9000/3 = 3000$ \\
        \hline
        $x_6$ & 0 & 8000 & 3 & 1 & 2 & 0 & 0 & 1 & $8000/2 = 4000$ \\
        \hline
        \multicolumn{2}{|c|}{$\Delta_j$} & \textbf{Z = 0} & -3 & -2 & \textbf{-5} $\uparrow$ & 0 & 0 & 0 & \\
        \hline
    \end{tabular}
\end{table}
\textit{Giải thích xoay bảng 1:}
\begin{itemize}[itemsep=0pt, parsep=0pt, topsep=0pt]
    \item Xét dòng $\Delta_j$, ta thấy $\Delta_3 = -5$ là số âm nhỏ nhất $\Rightarrow$ Chọn biến $x_3$ đưa \textbf{vào cơ sở}.
    \item Lập tỷ số $\lambda_i = PA / a_{i3}$ (với $a_{i3} > 0$), ta được $\min\{2000, 3000, 4000\} = 2000$ $\Rightarrow$ Chọn biến $x_4$ đưa \textbf{ra khỏi cơ sở}. Phần tử trụ là $\mathbf{2}$.
\end{itemize}

\vspace{0.3cm}
\textbf{Bảng 2 (Bước lặp 1):} Thực hiện phép biến đổi Gauss-Jordan theo phần tử trụ $2$.
\begin{table}[H]
    \centering
    \renewcommand{\arraystretch}{1.3}
    \begin{tabular}{|c|c|c|c|c|c|c|c|c|c|}
        \hline
        \multirow{2}{*}{\textbf{Cơ sở}} & \multirow{2}{*}{$\mathbf{C_b}$} & \multirow{2}{*}{\textbf{PA}} & $x_1$ & $x_2$ & $x_3$ & $x_4$ & $x_5$ & $x_6$ & \multirow{2}{*}{\textbf{Tỷ số $\lambda$}} \\
        \cline{4-9}
        & & & 3 & 2 & 5 & 0 & 0 & 0 & \\
        \hline
        $x_3$ & 5 & 2000 & 1/2 & 1/2 & 1 & 1/2 & 0 & 0 & $2000 / (1/2) = 4000$ \\
        \hline
        $x_5$ & 0 & 3000 & -1/2 & -1/2 & 0 & -3/2 & 1 & 0 & - \\
        \hline
        $x_6$ & 0 & 4000 & \textbf{2} & 0 & 0 & -1 & 0 & 1 & $4000 / 2 = \mathbf{2000} \rightarrow$ \\
        \hline
        \multicolumn{2}{|c|}{$\Delta_j$} & \textbf{Z = 10000} & \textbf{-1/2} $\uparrow$ & 1/2 & 0 & 5/2 & 0 & 0 & \\
        \hline
    \end{tabular}
\end{table}
\textit{Giải thích xoay bảng 2:}
\begin{itemize}[itemsep=0pt, parsep=0pt, topsep=0pt]
    \item Vẫn còn $\Delta_1 = -1/2 < 0$, do đó phương án hiện tại chưa tối ưu. Chọn $x_1$ đưa \textbf{vào cơ sở}.
    \item Lập tỷ số $\lambda$ ở cột $x_1$ (chỉ tính dòng có $a_{i1} > 0$), ta có $\min\{4000, 2000\} = 2000$ $\Rightarrow$ Chọn $x_6$ \textbf{ra khỏi cơ sở}. Phần tử trụ là $\mathbf{2}$.
\end{itemize}

\vspace{0.3cm}
\textbf{Bảng 3 (Bước lặp 2):} Cập nhật bảng theo phần tử trụ $2$.
\begin{table}[H]
    \centering
    \renewcommand{\arraystretch}{1.3}
    \begin{tabular}{|c|c|c|c|c|c|c|c|c|}
        \hline
        \multirow{2}{*}{\textbf{Cơ sở}} & \multirow{2}{*}{$\mathbf{C_b}$} & \multirow{2}{*}{\textbf{PA}} & $x_1$ & $x_2$ & $x_3$ & $x_4$ & $x_5$ & $x_6$ \\
        \cline{4-9}
        & & & 3 & 2 & 5 & 0 & 0 & 0 \\
        \hline
        $x_3$ & 5 & 1000 & 0 & 1/2 & 1 & 3/4 & 0 & -1/4 \\
        \hline
        $x_5$ & 0 & 4000 & 0 & -1/2 & 0 & -7/4 & 1 & 1/4 \\
        \hline
        $x_1$ & 3 & 2000 & 1 & 0 & 0 & -1/2 & 0 & 1/2 \\
        \hline
        \multicolumn{2}{|c|}{$\Delta_j$} & \textbf{Z = 11000} & 0 & 1/2 & 0 & 9/4 & 0 & 1/4 \\
        \hline
    \end{tabular}
\end{table}

\textbf{3. Kết luận}

Tại Bảng 3, tất cả các hệ số ước lượng đều thỏa mãn $\Delta_j \ge 0, \forall j$. Thuật toán dừng lại.
\begin{itemize}
    \item Phương án tối ưu của bài toán là: $x^* = (2000, 0, 1000)$.
    \item Giá trị lớn nhất của hàm mục tiêu là: $f(x)_{\max} = 11000$.
\end{itemize}